% Options for packages loaded elsewhere
\PassOptionsToPackage{unicode}{hyperref}
\PassOptionsToPackage{hyphens}{url}
\PassOptionsToPackage{dvipsnames,svgnames,x11names}{xcolor}
%
\documentclass[
]{article}
\usepackage{amsmath,amssymb}
\usepackage{lmodern}
\usepackage{iftex}
\ifPDFTeX
  \usepackage[T1]{fontenc}
  \usepackage[utf8]{inputenc}
  \usepackage{textcomp} % provide euro and other symbols
\else % if luatex or xetex
  \usepackage{unicode-math}
  \defaultfontfeatures{Scale=MatchLowercase}
  \defaultfontfeatures[\rmfamily]{Ligatures=TeX,Scale=1}
\fi
% Use upquote if available, for straight quotes in verbatim environments
\IfFileExists{upquote.sty}{\usepackage{upquote}}{}
\IfFileExists{microtype.sty}{% use microtype if available
  \usepackage[]{microtype}
  \UseMicrotypeSet[protrusion]{basicmath} % disable protrusion for tt fonts
}{}
\makeatletter
\@ifundefined{KOMAClassName}{% if non-KOMA class
  \IfFileExists{parskip.sty}{%
    \usepackage{parskip}
  }{% else
    \setlength{\parindent}{0pt}
    \setlength{\parskip}{6pt plus 2pt minus 1pt}}
}{% if KOMA class
  \KOMAoptions{parskip=half}}
\makeatother
\usepackage{xcolor}
\IfFileExists{xurl.sty}{\usepackage{xurl}}{} % add URL line breaks if available
\IfFileExists{bookmark.sty}{\usepackage{bookmark}}{\usepackage{hyperref}}
\hypersetup{
  colorlinks=true,
  linkcolor={blue},
  filecolor={Maroon},
  citecolor={Blue},
  urlcolor={blue},
  pdfcreator={LaTeX via pandoc}}
\urlstyle{same} % disable monospaced font for URLs
\usepackage[margin=1in]{geometry}
\usepackage{color}
\usepackage{fancyvrb}
\newcommand{\VerbBar}{|}
\newcommand{\VERB}{\Verb[commandchars=\\\{\}]}
\DefineVerbatimEnvironment{Highlighting}{Verbatim}{commandchars=\\\{\}}
% Add ',fontsize=\small' for more characters per line
\usepackage{framed}
\definecolor{shadecolor}{RGB}{248,248,248}
\newenvironment{Shaded}{\begin{snugshade}}{\end{snugshade}}
\newcommand{\AlertTok}[1]{\textcolor[rgb]{0.94,0.16,0.16}{#1}}
\newcommand{\AnnotationTok}[1]{\textcolor[rgb]{0.56,0.35,0.01}{\textbf{\textit{#1}}}}
\newcommand{\AttributeTok}[1]{\textcolor[rgb]{0.77,0.63,0.00}{#1}}
\newcommand{\BaseNTok}[1]{\textcolor[rgb]{0.00,0.00,0.81}{#1}}
\newcommand{\BuiltInTok}[1]{#1}
\newcommand{\CharTok}[1]{\textcolor[rgb]{0.31,0.60,0.02}{#1}}
\newcommand{\CommentTok}[1]{\textcolor[rgb]{0.56,0.35,0.01}{\textit{#1}}}
\newcommand{\CommentVarTok}[1]{\textcolor[rgb]{0.56,0.35,0.01}{\textbf{\textit{#1}}}}
\newcommand{\ConstantTok}[1]{\textcolor[rgb]{0.00,0.00,0.00}{#1}}
\newcommand{\ControlFlowTok}[1]{\textcolor[rgb]{0.13,0.29,0.53}{\textbf{#1}}}
\newcommand{\DataTypeTok}[1]{\textcolor[rgb]{0.13,0.29,0.53}{#1}}
\newcommand{\DecValTok}[1]{\textcolor[rgb]{0.00,0.00,0.81}{#1}}
\newcommand{\DocumentationTok}[1]{\textcolor[rgb]{0.56,0.35,0.01}{\textbf{\textit{#1}}}}
\newcommand{\ErrorTok}[1]{\textcolor[rgb]{0.64,0.00,0.00}{\textbf{#1}}}
\newcommand{\ExtensionTok}[1]{#1}
\newcommand{\FloatTok}[1]{\textcolor[rgb]{0.00,0.00,0.81}{#1}}
\newcommand{\FunctionTok}[1]{\textcolor[rgb]{0.00,0.00,0.00}{#1}}
\newcommand{\ImportTok}[1]{#1}
\newcommand{\InformationTok}[1]{\textcolor[rgb]{0.56,0.35,0.01}{\textbf{\textit{#1}}}}
\newcommand{\KeywordTok}[1]{\textcolor[rgb]{0.13,0.29,0.53}{\textbf{#1}}}
\newcommand{\NormalTok}[1]{#1}
\newcommand{\OperatorTok}[1]{\textcolor[rgb]{0.81,0.36,0.00}{\textbf{#1}}}
\newcommand{\OtherTok}[1]{\textcolor[rgb]{0.56,0.35,0.01}{#1}}
\newcommand{\PreprocessorTok}[1]{\textcolor[rgb]{0.56,0.35,0.01}{\textit{#1}}}
\newcommand{\RegionMarkerTok}[1]{#1}
\newcommand{\SpecialCharTok}[1]{\textcolor[rgb]{0.00,0.00,0.00}{#1}}
\newcommand{\SpecialStringTok}[1]{\textcolor[rgb]{0.31,0.60,0.02}{#1}}
\newcommand{\StringTok}[1]{\textcolor[rgb]{0.31,0.60,0.02}{#1}}
\newcommand{\VariableTok}[1]{\textcolor[rgb]{0.00,0.00,0.00}{#1}}
\newcommand{\VerbatimStringTok}[1]{\textcolor[rgb]{0.31,0.60,0.02}{#1}}
\newcommand{\WarningTok}[1]{\textcolor[rgb]{0.56,0.35,0.01}{\textbf{\textit{#1}}}}
\usepackage{longtable,booktabs,array}
\usepackage{calc} % for calculating minipage widths
% Correct order of tables after \paragraph or \subparagraph
\usepackage{etoolbox}
\makeatletter
\patchcmd\longtable{\par}{\if@noskipsec\mbox{}\fi\par}{}{}
\makeatother
% Allow footnotes in longtable head/foot
\IfFileExists{footnotehyper.sty}{\usepackage{footnotehyper}}{\usepackage{footnote}}
\makesavenoteenv{longtable}
\setlength{\emergencystretch}{3em} % prevent overfull lines
\providecommand{\tightlist}{%
  \setlength{\itemsep}{0pt}\setlength{\parskip}{0pt}}
\setcounter{secnumdepth}{-\maxdimen} % remove section numbering
\ifLuaTeX
  \usepackage{selnolig}  % disable illegal ligatures
\fi

\author{}
\date{}

\begin{document}

\hypertarget{ares-laptop-disaster-recovery-runbook}{%
\section{Ares Laptop --- Disaster Recovery
Runbook}\label{ares-laptop-disaster-recovery-runbook}}

\textbf{Tags:} \#runbook \#disaster-recovery \#ares \#restic \#ops
\textbf{Created:} 2026-07-20 \textbf{Scope:} Total loss of the Ares
admin/jump laptop → back to a working box.

The next hardware death should be fast and boring. This is the
do-this-next-time procedure, ordered start to finish. It was written
from the 2026-07 recovery of the old Ares (an \textasciitilde8-yr-old
laptop that died mid-session with a motherboard power-delivery failure)
onto a \textbf{Dell Latitude 3580} (service tag \texttt{5HWCNJ2}).

Full narrative + evidence:
\href{ares-recovery-context.md}{\texttt{ares-recovery-context.md}} (Part
1 = claude.ai chat reconstruction; Part 1B = machine-history ground
truth). The restic mechanism itself is also documented in
\href{../vault/Runbook/Ares-Backup-Restore.md}{\texttt{../vault/Runbook/Ares-Backup-Restore.md}}
--- this runbook is the wider ``hardware is dead, rebuild everything''
superset.

\begin{quote}
\textbf{Ares is only the admin box, not a cluster member.} When it dies
the cluster keeps running --- DNS, Proxmox, OPNsense, backups all stay
up. There is no production outage clock. Work the steps in order; don't
rush hardware.
\end{quote}

\textbf{Key facts to keep this runbook accurate (verbatim --- do not
paraphrase):}

\begin{longtable}[]{@{}
  >{\raggedright\arraybackslash}p{(\columnwidth - 2\tabcolsep) * \real{0.5000}}
  >{\raggedright\arraybackslash}p{(\columnwidth - 2\tabcolsep) * \real{0.5000}}@{}}
\toprule
\begin{minipage}[b]{\linewidth}\raggedright
Item
\end{minipage} & \begin{minipage}[b]{\linewidth}\raggedright
Value
\end{minipage} \\
\midrule
\endhead
Backup tool & \texttt{restic} (NOT PBS --- see §5) \\
Repo & \texttt{sftp:randy:/mnt/bulk/backups/ares} (encrypted) \\
Backup host & Randy --- \texttt{192.168.10.187}, \texttt{bulk} ZFS
pool \\
\texttt{randy} SSH alias & \texttt{root@192.168.10.187} (in
\texttt{\textasciitilde{}/.ssh/config}) --- resolves to \textbf{root},
which is why SFTP can read the repo \\
Repo dir perms & root-owned, mode \texttt{700} --- only root can read
it \\
Passphrase (on-disk) &
\texttt{\textasciitilde{}/.config/restic/ares-randy.pass} (mode 600) \\
Passphrase (offsite) & Vaultwarden entry
\texttt{Ares\ restic\ backup\ -\textgreater{}\ Randy\ (repo\ password)}
(pve3 LXC 102) \\
Hostname & \texttt{ares} --- \textbf{reuse it} (see §2) \\
Ares IP & \texttt{192.168.10.199} (mgmt VLAN 1) \\
\bottomrule
\end{longtable}

\begin{center}\rule{0.5\linewidth}{0.5pt}\end{center}

\hypertarget{pre-disaster-checklist-the-prevention-layer}{%
\subsection{§0 --- Pre-disaster checklist (the prevention
layer)}\label{pre-disaster-checklist-the-prevention-layer}}

Do these while Ares is healthy. Each one directly collapses a step that
cost real time in the last recovery. This is the highest-leverage
section in the document.

\begin{itemize}
\tightlist
\item[$\square$]
  \textbf{Offline copy of the three secrets that gate everything else.}
  Put the restic passphrase \textbf{and} the SSH private keys
  (\texttt{\textasciitilde{}/.ssh/id\_ed25519*}) \textbf{and} the
  Ansible vault key somewhere that survives the laptop dying: a phone
  password manager, a printed copy in a safe, or the salvaged NVMe in an
  enclosure (below). \emph{Why:} last time the passphrase lived only in
  Vaultwarden, and reaching it needed a from-scratch nginx TLS proxy on
  a machine with nothing installed (§4). An offline copy collapses §4 to
  a copy-paste.
\item[$\square$]
  \textbf{Keep a USB 3.2 NVMe M-key enclosure in the kit permanently}
  (ORICO / UGREEN, \textasciitilde\$15--20). \emph{Why:} it is the
  fastest path to the old drive's secrets and the ultimate fallback
  source of truth. The salvaged 970 EVO Plus is currently unreadable
  precisely because there's no enclosure on hand.
\item[$\square$]
  \textbf{The backup mechanism is documented and unambiguous:} restic at
  \texttt{sftp:randy:/mnt/bulk/backups/ares}, repo dir is
  \textbf{root-owned mode 700}, so restic must run \textbf{as root} (the
  \texttt{randy} alias already is root). It is \textbf{not PBS}. See §5.
  \emph{Why:} last time recovery started down the PBS /
  \texttt{proxmox-backup-client} path and lost time before discovering
  it was restic.
\item[$\square$]
  \textbf{A working Debian DVD-1 ISO is on hand} (not just netinst) with
  a way to flash it (Rufus). \emph{Why:} the Intel Wireless-AC card
  needs non-free firmware the netinst can't fetch without a network
  (§3).
\item[$\square$]
  \textbf{The backup is monitored and freshness is trusted.} The
  \texttt{AresBackupStale} Grafana rule (\textgreater26h → Discord) must
  be live and believed. \emph{Why:} verify snapshot age FIRST every time
  (§1); don't assume the backup is current.
\end{itemize}

\begin{center}\rule{0.5\linewidth}{0.5pt}\end{center}

\hypertarget{confirm-the-backup-is-good-before-touching-hardware}{%
\subsection{§1 --- Confirm the backup is good BEFORE touching
hardware}\label{confirm-the-backup-is-good-before-touching-hardware}}

Prove you have a restorable, recent, intact snapshot \textbf{before}
buying or wiping anything. If this step fails, fix the backup story
first --- everything downstream assumes it.

From any machine that can reach Randy and has the passphrase (this can
be done before the new laptop even exists --- e.g.~from another box or a
phone SSH client):

\begin{Shaded}
\begin{Highlighting}[]
\FunctionTok{ssh}\NormalTok{ randy }\StringTok{\textquotesingle{}restic {-}r /mnt/bulk/backups/ares snapshots\textquotesingle{}}
\end{Highlighting}
\end{Shaded}

Lists every snapshot with ID + timestamp. Run \textbf{on Randy as root}
(the \texttt{randy} alias is root, and the repo is root-only), so you
avoid the SFTP permission trap entirely for a read-only check. Confirm
the newest snapshot is recent --- the last real recovery found the
newest at \textbf{2026-07-17 04:03, 110.212 GiB}.

\begin{Shaded}
\begin{Highlighting}[]
\FunctionTok{ssh}\NormalTok{ randy }\StringTok{\textquotesingle{}restic {-}r /mnt/bulk/backups/ares check\textquotesingle{}}
\end{Highlighting}
\end{Shaded}

Verifies repo structure and metadata are intact (fast). This is
\emph{not} a full data re-read --- for that,
\texttt{restic\ check\ -\/-read-data} pulls back all \textasciitilde110
GiB; skip it routinely (ZFS checksums on \texttt{bulk} already guard
bit-rot) but consider it once before you delete any staging copy.

\begin{quote}
\textbf{Trap from last time:} don't judge freshness by reading the
backup \emph{log} --- and especially not the log \emph{inside a restored
copy}, which can be a mid-run snapshot that ends at ``backup starting''
with no completion line. Trust \texttt{restic\ snapshots} against the
\textbf{live} repo, not a log artifact.
\end{quote}

\begin{quote}
\textbf{If a stale lock blocks you:} \texttt{restic\ unlock} clears an
abandoned lock left by the dead machine.
\end{quote}

\begin{center}\rule{0.5\linewidth}{0.5pt}\end{center}

\hypertarget{replacement-hardware-selection}{%
\subsection{§2 --- Replacement hardware
selection}\label{replacement-hardware-selection}}

Check interfaces and RAM \textbf{before you buy or commit a drive} ---
the last stopgap laptop had two nasty surprises that a spec check would
have caught.

\textbf{Interface / RAM checklist (verify against the actual model's
service manual):}

\begin{itemize}
\tightlist
\item[$\square$]
  \textbf{Drive interface: NVMe vs SATA, and is there an M.2 NVMe slot
  at all?} The Latitude 3580 has \textbf{no M.2 NVMe slot} --- only a
  2.5'' SATA bay + WLAN/WWAN. Its RWMDF caddy carries \textbf{M.2 SATA
  only}, not NVMe/PCIe. The salvaged 1TB 970 EVO Plus (NVMe) therefore
  \textbf{cannot be used internally} --- a WD Black 2.5'' SATA drive
  went in the bay instead. Confirm the slot type before planning to
  reuse any drive.
\item[$\square$]
  \textbf{RAM is socketed, not soldered.} (A ThinkPad T14 Gen 2 was
  rejected for soldered RAM.) Socketed = upgradable and repairable.
\item[$\square$]
  Enough of everything for a jump box: the old Ares had 32GB RAM.
\end{itemize}

\textbf{Reuse the hostname \texttt{ares}.} This is a rule, not a
preference. Keeping the name means you do \textbf{not} have to touch:

\begin{itemize}
\tightlist
\item
  Ansible inventory
\item
  Pi-hole DNS records
\item
  DHCP MAC reservations
\item
  Headscale node entries
\item
  every script that hardcodes \texttt{ares}
\end{itemize}

On any machine that had SSH'd to the \emph{old} Ares, clear its stale
host key so you don't hit ``host key verification failed'' against the
new box:

\begin{Shaded}
\begin{Highlighting}[]
\FunctionTok{ssh{-}keygen} \AttributeTok{{-}R}\NormalTok{ ares}
\FunctionTok{ssh{-}keygen} \AttributeTok{{-}R}\NormalTok{ 192.168.10.199}
\end{Highlighting}
\end{Shaded}

Removes the old fingerprint for the name and the IP from
\texttt{known\_hosts}. (This exact class of stale-key failure once
masqueraded as an ``auth failure'' against the EX3400 --- see
\texttt{EX3400-SSH-Auth-Failure-RCA.md}.)

\begin{center}\rule{0.5\linewidth}{0.5pt}\end{center}

\hypertarget{install-debian}{%
\subsection{§3 --- Install Debian}\label{install-debian}}

Target: \textbf{Debian 13 Trixie + KDE Plasma 6} (better 7th-gen Intel
firmware support than Debian 12).

\textbf{Media:} - Use the full \textbf{DVD-1 ISO}, \emph{not} netinst.
\emph{Why:} the Intel Wireless-AC card needs non-free firmware; netinst
assumes a working network to fetch it, which you don't have yet. The DVD
carries firmware + the full KDE set offline. - Flash with \textbf{Rufus
in DD mode} (not ISO mode). \emph{Why:} ISO mode rewrites the boot
sector and corrupts Debian's hybrid ISO; DD writes it byte-for-byte.
Target \textbf{GPT / UEFI}.

\textbf{Installer:} - At software selection, \textbf{UNCHECK GNOME,
CHECK KDE Plasma} --- Debian defaults to GNOME and it's easy to blow
past. - Set hostname \textbf{\texttt{ares}} (see §2).

\hypertarget{uefi-secure-boot-grub-nvram-recovery-expect-this-on-dell}{%
\subsubsection{§3.1 --- UEFI / Secure-Boot / GRUB-NVRAM recovery (expect
this on
Dell)}\label{uefi-secure-boot-grub-nvram-recovery-expect-this-on-dell}}

GRUB installs fine but the machine boots to firmware with no
\texttt{debian} entry --- the EFI files are on the ESP but the NVRAM
boot variable didn't stick (it clears when the USB is pulled and the
firmware reshuffles boot order). Two fixes; the first is what worked.

\textbf{Fix A --- add the boot entry by hand in firmware (what worked):}
- In the Dell firmware boot menu, add a boot option pointing at
\texttt{\textbackslash{}EFI\textbackslash{}debian\textbackslash{}shimx64.efi}
(fallback
\texttt{\textbackslash{}EFI\textbackslash{}debian\textbackslash{}grubx64.efi}),
name it \texttt{debian}, and move it to the \textbf{top} of the boot
order. - \textbf{Secure Boot:} Dell firmware can silently re-enable it
after changes, and then shim refuses to load. Re-confirm \textbf{Secure
Boot OFF} under \textbf{F2 → Boot Configuration}.

\textbf{Fix B --- rescue mode (same install USB → Advanced → Rescue):}

\begin{Shaded}
\begin{Highlighting}[]
\ExtensionTok{grub{-}install} \AttributeTok{{-}{-}target}\OperatorTok{=}\NormalTok{x86\_64{-}efi }\AttributeTok{{-}{-}efi{-}directory}\OperatorTok{=}\NormalTok{/boot/efi }\AttributeTok{{-}{-}bootloader{-}id}\OperatorTok{=}\NormalTok{debian }\AttributeTok{{-}{-}recheck}
\ExtensionTok{update{-}grub}
\ExtensionTok{efibootmgr} \AttributeTok{{-}v}
\end{Highlighting}
\end{Shaded}

Reinstalls GRUB's EFI files, rewrites the NVRAM boot entry named
\texttt{debian}, and prints the boot order so you can confirm it took.
Remove leftover Windows entries with:

\begin{Shaded}
\begin{Highlighting}[]
\ExtensionTok{efibootmgr} \AttributeTok{{-}b} \OperatorTok{\textless{}}\NormalTok{NNNN}\OperatorTok{\textgreater{}}\NormalTok{ {-}B}
\end{Highlighting}
\end{Shaded}

Deletes boot entry number \texttt{NNNN} (read the number from
\texttt{efibootmgr\ -v} output).

\begin{center}\rule{0.5\linewidth}{0.5pt}\end{center}

\hypertarget{retrieve-the-restic-passphrase}{%
\subsection{§4 --- Retrieve the restic
passphrase}\label{retrieve-the-restic-passphrase}}

You need the repo passphrase before you can restore (§5). Take the fast
path if you prepared §0; the Vaultwarden path is the fallback only.

\hypertarget{fast-path-use-this-if-0-was-done}{%
\subsubsection{§4.1 --- Fast path (use this if §0 was
done)}\label{fast-path-use-this-if-0-was-done}}

\begin{itemize}
\tightlist
\item
  \textbf{Offline copy:} phone password manager, printed copy, or
  wherever you stashed it. Copy it and go straight to §5.
\item
  \textbf{Salvaged NVMe via enclosure:} put the old 970 EVO Plus in the
  USB NVMe enclosure, mount it read-only, and read the passphrase
  directly from
  \texttt{\textasciitilde{}/.config/restic/ares-randy.pass} (the SSH
  keys and Ansible vault key are right there too). This makes the whole
  Vaultwarden detour unnecessary.
\end{itemize}

\hypertarget{fallback-pull-it-from-vaultwarden}{%
\subsubsection{§4.2 --- Fallback: pull it from
Vaultwarden}\label{fallback-pull-it-from-vaultwarden}}

Vaultwarden runs on pve3 (LXC 102) and survives Ares dying, but it needs
HTTPS and can't be hit directly from a fresh laptop. Two ways in:

\textbf{Option 1 --- read the secret straight out of the container over
SSH} (fastest; skips the browser entirely). If the entry you need is
stored as a container env var / file, \texttt{pct\ exec} into it from a
Proxmox node:

\begin{Shaded}
\begin{Highlighting}[]
\FunctionTok{ssh}\NormalTok{ root@192.168.10.201 }\StringTok{"pct exec 102 {-}{-} grep ADMIN\_TOKEN /etc/vaultwarden.env"}
\end{Highlighting}
\end{Shaded}

\texttt{192.168.10.201} = pve3, \texttt{102} = the Vaultwarden LXC.
(This exact command pulled the Vaultwarden admin token during the last
recovery. The restic passphrase is a \emph{vault item}, not an env var,
so this gets you the admin token / into the box --- you may still need
the UI or \texttt{bw} CLI to read the item itself.)

\textbf{Option 2 --- stand up a throwaway TLS proxy and use the web UI}
(what was done last time). Vaultwarden speaks plain HTTP on
\texttt{:8080}; the browser wants HTTPS. Tunnel to it, then front it
with a self-signed TLS proxy:

\begin{Shaded}
\begin{Highlighting}[]
\ExtensionTok{ss} \AttributeTok{{-}tlnp} \KeywordTok{|} \FunctionTok{grep} \AttributeTok{{-}E} \StringTok{\textquotesingle{}8080|8443\textquotesingle{}}          \CommentTok{\# confirm the tunnel/listener is up}
\FunctionTok{sudo}\NormalTok{ apt install }\AttributeTok{{-}y}\NormalTok{ nginx openssl        }\CommentTok{\# neither is on a fresh box}
\FunctionTok{sudo}\NormalTok{ openssl req }\AttributeTok{{-}x509} \AttributeTok{{-}newkey}\NormalTok{ rsa:2048 }\AttributeTok{{-}nodes} \AttributeTok{{-}days}\NormalTok{ 30 }\DataTypeTok{\textbackslash{}}
  \AttributeTok{{-}keyout}\NormalTok{ /etc/ssl/private/vw.key }\AttributeTok{{-}out}\NormalTok{ /etc/ssl/certs/vw.crt }\AttributeTok{{-}subj} \StringTok{"/CN=localhost"}
\end{Highlighting}
\end{Shaded}

First line verifies the SSH local-forward to \texttt{:8080} is
listening. The \texttt{openssl} line mints a 30-day self-signed cert for
\texttt{localhost}. Then add an nginx server block listening
\texttt{8443\ ssl} that proxies to \texttt{127.0.0.1:8080}, validate and
load it:

\begin{Shaded}
\begin{Highlighting}[]
\FunctionTok{sudo}\NormalTok{ nginx }\AttributeTok{{-}t} \KeywordTok{\&\&} \FunctionTok{sudo}\NormalTok{ systemctl restart nginx}
\end{Highlighting}
\end{Shaded}

\texttt{nginx\ -t} checks the config before restart. Open
\texttt{https://localhost:8443} in Firefox, accept the self-signed cert,
log in, and copy the entry
\textbf{\texttt{Ares\ restic\ backup\ -\textgreater{}\ Randy\ (repo\ password)}}.

\begin{quote}
\textbf{After recovery, close this gap for good (§0):} get the
passphrase into an offline store so §4.2 never happens again. During the
last recovery the restic entry was \emph{added} to Vaultwarden as
remediation --- also keep a copy off-cluster.
\end{quote}

\begin{center}\rule{0.5\linewidth}{0.5pt}\end{center}

\hypertarget{restore-with-restic}{%
\subsection{§5 --- Restore with restic}\label{restore-with-restic}}

This is \textbf{restic}, not PBS. Do \textbf{not} reach for
\texttt{proxmox-backup-client} / \texttt{.pxar} /
\texttt{host/ares/\textless{}timestamp\textgreater{}} --- that path was
a dead end last time.

\textbf{Repo reality:} \texttt{/mnt/bulk/backups/ares} on Randy is
\textbf{root-owned, mode 700}. Only root can read it. The trick that
makes this painless: the \texttt{randy} SSH alias in
\texttt{\textasciitilde{}/.ssh/config} already points at
\textbf{\texttt{root@192.168.10.187}}, so \texttt{sftp:randy:} is an
SFTP session \emph{as root} and can read the repo --- you do
\textbf{not} have to log into Randy and run restic there. (For a
read-only \texttt{snapshots}/\texttt{check} it's still fine to run on
Randy directly, as in §1.)

\textbf{Step 1 --- get restic on the fresh box} (it won't be installed):

\begin{Shaded}
\begin{Highlighting}[]
\FunctionTok{sudo}\NormalTok{ apt install }\AttributeTok{{-}y}\NormalTok{ restic}
\end{Highlighting}
\end{Shaded}

Installs the restic binary. (The healthy Ares also keeps a static binary
at \texttt{\textasciitilde{}/.local/bin/restic}; either works.)

\textbf{Step 2 --- point restic at the repo and unlock it:}

\begin{Shaded}
\begin{Highlighting}[]
\BuiltInTok{export} \VariableTok{RESTIC\_REPOSITORY}\OperatorTok{=}\StringTok{"sftp:randy:/mnt/bulk/backups/ares"}
\BuiltInTok{export} \VariableTok{RESTIC\_PASSWORD\_FILE}\OperatorTok{=}\NormalTok{\textasciitilde{}/.config/restic/ares{-}randy.pass   }\CommentTok{\# if you restored the file}
\CommentTok{\# — or, with the passphrase from §4, avoid putting it in shell history:}
\CommentTok{\# read {-}rs RESTIC\_PASSWORD; export RESTIC\_PASSWORD}
\ExtensionTok{restic}\NormalTok{ snapshots}
\end{Highlighting}
\end{Shaded}

\texttt{RESTIC\_REPOSITORY} selects the repo over SFTP-as-root; the
password is supplied by file (preferred) or env.
\texttt{restic\ snapshots} confirms you can actually read it and shows
the snapshot you're about to restore.

\begin{quote}
\textbf{Security note (learned the hard way):} last time the passphrase
was inlined as
\texttt{export\ RESTIC\_PASSWORD=\textquotesingle{}…\textquotesingle{}},
which left it in \textbf{plaintext in
\texttt{\textasciitilde{}/.bash\_history}}. Prefer
\texttt{RESTIC\_PASSWORD\_FILE}, or \texttt{read\ -rs} so it never hits
history. If you do inline it, scrub the line from
\texttt{\textasciitilde{}/.bash\_history} afterward.
\end{quote}

\textbf{Step 3 --- restore to staging, NEVER straight onto live
\texttt{/home}:}

\begin{Shaded}
\begin{Highlighting}[]
\ExtensionTok{restic}\NormalTok{ restore latest }\AttributeTok{{-}{-}target}\NormalTok{ \textasciitilde{}/restore}
\end{Highlighting}
\end{Shaded}

Restores the newest snapshot into \texttt{\textasciitilde{}/restore} on
the laptop (not onto live \texttt{/home}, which KDE has dotfiles open in
--- overwriting live is unpredictable). This runs from the laptop over
SFTP-as-root and lands directly in the staging dir; no Randy-side
staging or rsync is needed. To pull just one tree, add
\texttt{-\/-include\ /home/machismo/Home-Lab}.

\textbf{Step 4 --- inspect, then selectively reinstate.} With the old
home sitting under \texttt{\textasciitilde{}/restore/home/machismo/},
copy back what you need --- dotfiles, \texttt{\textasciitilde{}/.ssh/}
keys (check perms: \texttt{600} on private keys), the \texttt{Home-Lab}
repo, \texttt{\textasciitilde{}/.config/restic/}, and the backup
script/timer. Do it deliberately, not with a blind overwrite.

\textbf{Step 5 --- repopulate excluded content.} Upstream git submodules
(\texttt{claude-desktop-debian}, \texttt{pacextractor},
\texttt{spreadtrum\_flash}, \texttt{CVE-2022-38694\_unlock\_bootloader})
and caches/venvs were excluded from the backup to save space:

\begin{Shaded}
\begin{Highlighting}[]
\FunctionTok{git} \AttributeTok{{-}C}\NormalTok{ \textasciitilde{}/dotfiles submodule update }\AttributeTok{{-}{-}init} \AttributeTok{{-}{-}recursive}
\end{Highlighting}
\end{Shaded}

Re-clones the excluded submodules from their own remotes. Rebuild Python
venvs from their \texttt{requirements.txt} as needed.

\textbf{Step 6 --- reinstate the backup tooling immediately} (this is
what stopped backups last time --- the rebuild wiped restic,
\texttt{ares-backup.sh}, \texttt{\textasciitilde{}/.config/restic/}, and
all user timers, so \texttt{AresBackupStale} fired). Restore the
\texttt{restic} binary to \texttt{\textasciitilde{}/.local/bin/},
\texttt{ares-backup.sh}, and \texttt{\textasciitilde{}/.config/restic/},
then enable the timer and run once:

\begin{Shaded}
\begin{Highlighting}[]
\ExtensionTok{systemctl} \AttributeTok{{-}{-}user}\NormalTok{ enable }\AttributeTok{{-}{-}now}\NormalTok{ ares{-}backup.timer}
\ExtensionTok{systemctl} \AttributeTok{{-}{-}user}\NormalTok{ list{-}timers ares{-}backup.timer}
\end{Highlighting}
\end{Shaded}

Enables the daily 04:00 backup timer and shows its next run, so the new
box is protected again the same night. Full detail:
\href{../vault/Runbook/Ares-Backup-Restore.md}{\texttt{../vault/Runbook/Ares-Backup-Restore.md}}.

\begin{quote}
\textbf{Clean up staging once reinstated:}
\texttt{\textasciitilde{}/restore} is temporarily excluded in
\texttt{\textasciitilde{}/.config/restic/excludes.txt} so it isn't
double-captured. Remove it (and the exclude line) once you're confident
everything is reinstated.
\end{quote}

\begin{center}\rule{0.5\linewidth}{0.5pt}\end{center}

\hypertarget{rebuild-the-toolchain-claude-code}{%
\subsection{§6 --- Rebuild the toolchain (Claude
Code)}\label{rebuild-the-toolchain-claude-code}}

Install Claude Code via the official installer:

\begin{Shaded}
\begin{Highlighting}[]
\ExtensionTok{curl} \AttributeTok{{-}fsSL}\NormalTok{ https://claude.ai/install.sh }\KeywordTok{|} \FunctionTok{bash}
\end{Highlighting}
\end{Shaded}

Installs the \texttt{claude} CLI into
\texttt{\textasciitilde{}/.local/bin} (last recovery landed v2.1.215).

\textbf{\texttt{claude:\ command\ not\ found} --- the PATH fix.}
\texttt{\textasciitilde{}/.local/bin} isn't on PATH for interactive
non-login shells: the export lived only in
\texttt{\textasciitilde{}/.profile} (login shells) and was missing from
\texttt{\textasciitilde{}/.bashrc} (the normal desktop-terminal case).
Add it to \texttt{\textasciitilde{}/.bashrc}:

\begin{Shaded}
\begin{Highlighting}[]
\FunctionTok{grep} \AttributeTok{{-}n} \StringTok{\textquotesingle{}local/bin\textquotesingle{}}\NormalTok{ \textasciitilde{}/.bashrc \textasciitilde{}/.profile                       }\CommentTok{\# see what\textquotesingle{}s already there}
\BuiltInTok{printf} \StringTok{\textquotesingle{}\textbackslash{}nexport PATH="$HOME/.local/bin:$PATH"\textbackslash{}n\textquotesingle{}} \OperatorTok{\textgreater{}\textgreater{}}\NormalTok{ \textasciitilde{}/.bashrc  }\CommentTok{\# append the export}
\BuiltInTok{source}\NormalTok{ \textasciitilde{}/.bashrc }\KeywordTok{\&\&} \FunctionTok{which}\NormalTok{ claude }\KeywordTok{\&\&} \ExtensionTok{claude} \AttributeTok{{-}{-}version}           \CommentTok{\# confirm it resolves}
\end{Highlighting}
\end{Shaded}

\texttt{grep} shows whether the export exists; \texttt{printf} appends
it cleanly (preferred over \texttt{echo} for predictable newlines);
\texttt{source} reloads the shell and the last two commands confirm
\texttt{claude} is now found.

\begin{quote}
\textbf{OPEN ISSUE --- Claude account (runbook can't fix this, only flag
it):} on the last rebuild, Claude login landed in an \textbf{empty
account with a payment prompt despite an active subscription} --- likely
a duplicate account from a different auth method. The account ID is
\textbf{not} usable as a CLI credential. Resolve it with Anthropic
support at \textbf{support.claude.com}; there is no local workaround.
\end{quote}

\begin{center}\rule{0.5\linewidth}{0.5pt}\end{center}

\hypertarget{post-restore-verification-checklist}{%
\subsection{§7 --- Post-restore verification
checklist}\label{post-restore-verification-checklist}}

Prove the rebuilt Ares is fully back in its admin/jump role. Work top to
bottom.

\begin{itemize}
\tightlist
\item[$\square$]
  \textbf{SSH keys intact:} \texttt{ls\ -l\ \textasciitilde{}/.ssh} ---
  private keys mode \texttt{600}; \texttt{randy}, \texttt{quarkylab},
  etc. aliases present in \texttt{\textasciitilde{}/.ssh/config}.
\item[$\square$]
  \textbf{Reach Randy as root:}
  \texttt{ssh\ randy\ \textquotesingle{}hostname\ \&\&\ zpool\ list\textquotesingle{}}
  --- confirms the alias, the key, and Randy's pools (\texttt{bulk},
  \texttt{datastore}) are healthy.
\item[$\square$]
  \textbf{SSH into each cluster node:} pve2 \texttt{.204}, pve3
  \texttt{.201}, pve4 \texttt{.202}, pve5 \texttt{.203}, plus pve1
  \texttt{.193}, QuarkyLab \texttt{.179}, Jarvis \texttt{.31}. Each
  should accept the key without a password.
\item[$\square$]
  \textbf{Proxmox reachable:} open a node UI,
  e.g.~\texttt{https://192.168.10.201:8006}.
\item[$\square$]
  \textbf{DNS resolves via Pi-hole:}
  \texttt{dig\ @192.168.10.177\ kylemason.org\ +short} (primary) and
  \texttt{@192.168.10.178} (secondary) both answer.
\item[$\square$]
  \textbf{Ansible sees the fleet (dry run, no changes):}
  \texttt{ansible\ all\ -m\ ping} then
  \texttt{ansible-playbook\ \textless{}play\textgreater{}\ -\/-check\ -\/-diff}.
  Confirms the inventory, the vault key, and SSH all work end-to-end.
\item[$\square$]
  \textbf{Backups running again:}
  \texttt{systemctl\ -\/-user\ list-timers\ ares-backup.timer} shows a
  next run; after the first run, \texttt{restic\ snapshots} shows a
  fresh-dated snapshot and \texttt{AresBackupStale} clears in Discord.
\item[$\square$]
  \textbf{git working trees clean:}
  \texttt{git\ -C\ \textasciitilde{}/dotfiles\ status} and
  \texttt{git\ -C\ \textasciitilde{}/Home-Lab\ status} --- restored, on
  the right branch, submodules initialised.
\item[$\square$]
  \textbf{Secrets hygiene:} confirm the restic passphrase is in an
  offline store (§0) and scrub any plaintext secret from
  \texttt{\textasciitilde{}/.bash\_history} (§5 security note).
\end{itemize}

\begin{center}\rule{0.5\linewidth}{0.5pt}\end{center}

\hypertarget{appendix-lessons-that-shaped-this-runbook}{%
\subsection{Appendix --- Lessons that shaped this
runbook}\label{appendix-lessons-that-shaped-this-runbook}}

\begin{enumerate}
\def\labelenumi{\arabic{enumi}.}
\tightlist
\item
  \textbf{Reuse the hostname} --- it kept
  Ansible/DNS/DHCP/Headscale/scripts untouched.
\item
  \textbf{Keep the secrets offline} --- the single biggest avoidable
  cost was the Vaultwarden-only passphrase and the from-scratch TLS
  proxy to reach it.
\item
  \textbf{Check the drive interface before buying/reusing} --- NVMe vs
  SATA, M.2 slot presence. The salvaged NVMe was dead-on-arrival for a
  SATA-only 3580.
\item
  \textbf{Keep a USB NVMe enclosure in the kit} --- the fallback source
  of truth.
\item
  \textbf{DVD ISO + Rufus DD} for offline-firmware laptops, and expect
  the Dell UEFI NVRAM / Secure-Boot boot-entry dance.
\item
  \textbf{It's restic, not PBS} ---
  \texttt{sftp:randy:/mnt/bulk/backups/ares}, root-owned 700, run as
  root. Enumerate ZFS pools before hunting for a repo.
\item
  \textbf{Verify snapshot freshness first, every time} --- against the
  live repo, not a log artifact.
\item
  \textbf{Restore to staging, inspect, then reinstate} --- never
  straight onto live \texttt{/home}.
\end{enumerate}

\end{document}
